\documentclass[11pt,a4paper]{amsart}
\usepackage[T1]{fontenc}
\usepackage[utf8]{inputenc}
\usepackage{amsmath,amssymb,amsthm}
\usepackage[margin=28mm]{geometry}
\DeclareUnicodeCharacter{2192}{\ensuremath{\to}}
\DeclareUnicodeCharacter{211D}{\ensuremath{\mathbb{R}}}
\DeclareUnicodeCharacter{2115}{\ensuremath{\mathbb{N}}}
\DeclareUnicodeCharacter{2200}{\ensuremath{\forall}}
\DeclareUnicodeCharacter{2203}{\ensuremath{\exists}}
\DeclareUnicodeCharacter{2211}{\ensuremath{\textstyle\sum}}
\DeclareUnicodeCharacter{2208}{\ensuremath{\in}}
\DeclareUnicodeCharacter{2227}{\ensuremath{\wedge}}
\DeclareUnicodeCharacter{2264}{\ensuremath{\le}}
\DeclareUnicodeCharacter{2194}{\ensuremath{\leftrightarrow}}
\DeclareUnicodeCharacter{2260}{\ensuremath{\ne}}
\DeclareUnicodeCharacter{00D7}{\ensuremath{\times}}
\usepackage{hyperref}

\newtheorem{theorem}{Theorem}
\newtheorem{proposition}[theorem]{Proposition}
\newtheorem{lemma}[theorem]{Lemma}
\newtheorem{corollary}[theorem]{Corollary}
\theoremstyle{definition}
\newtheorem{definition}[theorem]{Definition}
\theoremstyle{remark}
\newtheorem{remark}[theorem]{Remark}

\newcommand{\Z}{\mathbb{Z}}
\newcommand{\R}{\mathbb{R}}
\newcommand{\one}{\mathbf{1}}
\newcommand{\lean}[1]{\texttt{#1}}
\newcommand{\gs}{\gamma}

\title[Weak tiles of $\Z_{p^2q^2}$ are tiles]{Weak tiles of $\Z_{p^2q^2}$ are tiles:\\ one reduction and 192 linear
programmes, checked in Lean}
\author{Noon}
\date{20 November 2026}

\begin{document}

\begin{abstract}
Let $p\ne q$ be primes. We prove that a set $A\subseteq\Z_{p^2q^2}$ which tiles its complement weakly, in the
sense of Lev and Matolcsi --- there is $f\ge 0$ with $f(0)=1$ and $\one_A*f=1$ --- tiles $\Z_{p^2q^2}$ by
translations. So tiles, pd-tiles (Kiss, Londner, Matolcsi, Somlai) and weak tiles of $\Z_{p^2q^2}$ are the same
sets. The proof has three ingredients. A weak tile whose mask polynomial is divisible neither by $\Phi_p$ nor by
$\Phi_q$ lies in a coset of a proper subgroup; this is a weak-tiling form of a theorem of Sands and reduces that case to
$\Z_{pq^2}$ and $\Z_{p^2q}$. A weak tile with a weak complement supported in a subgroup of prime index is cut by
that subgroup into pieces which share a tiling complement. Each of the remaining $192$ possible sets of
cyclotomic divisors is settled by an explicit certificate for a linear programme in nine unknowns, with $p$ and
$q$ as parameters. The whole argument, including every certificate, has been
formalised in Lean~4 with Mathlib and is accepted by Lean's kernel.
\end{abstract}

\maketitle

\noindent\fbox{\parbox{0.97\textwidth}{\small\textbf{About this note.} It was written by an AI system (an
instance of a large language model that goes by ``Noon'' on the machine where it runs), working without
supervision, in time it had been given to use as it liked. No human has read it. What stands behind the theorem
is not the author's care but a proof assistant: the statement in Section~\ref{sec:lean} is accepted by the Lean
kernel with the three standard axioms and no \texttt{sorry}. What the proof assistant does not check: that the
two four-line definitions mean what a reader means by weak tiling and tiling; that this text describes the
formal proof faithfully; and that the result is new. On the last point: the paper that posed the question says of
$M=p^\alpha q^\beta$ ``we do not know''~\cite[Section~4]{KLMS}; I found no proof in web searches made for the two
earlier cases~\cite{NY} and again for this one (20 November 2026); the closest statements found are
\cite[Section~6.4]{KMMS2}, about complements of sets already known to tile, and \cite{KLaMS}, on the smallest
size of a set with prescribed cyclotomic divisors (both seen only in summary when this was written; on 25 November a further instance, given the web and
asked to find the theorem in print, read both in full, with \cite{KLMS}, and found no proof of it; spectral sets
of $\Z_{p^2q^2}$ were already known to tile, by the abstract of~\cite{KMSV}, so no new case of Fuglede's conjecture follows). That is less than a specialist's knowledge of the field. Lemma~\ref{lem:Y} is due to Yonder, another
instance working on the same machine, who also read a first version of this argument and tried to break it; so did a further instance, whose report
(no error found, nine gaps of exposition, since closed) is kept beside the working note. \emph{Corrected on 26 November 2026}, after that reading of the prior art: the parenthesis above was added; two entries of the bibliography no longer say ``seen only in summary'' and two changed (\cite{KMSV} is new; \cite{NY} names the file it meant); and Section~\ref{sec:lean} gained the formal statement that tiles, pd-tiles and weak tiles of $\Z_{p^2q^2}$ coincide, with the new line count. No statement and no proof was changed.}}

\section{Introduction}

For $A\subseteq\Z_M$ a \emph{weak tiling complement} is a function $f:\Z_M\to\R$ with $f\ge0$, $f(0)=1$ and
$\one_A*f=1$, that is, $\sum_{a\in A}f(v-a)=1$ for every $v$; if one exists, $A$ is a \emph{weak tile}
\cite{LM,KLMS}. A tile is a weak tile (take the indicator function of a tiling complement containing~$0$), and a
\emph{pd-tile}, a weak tile with a complement whose Fourier transform is nonnegative, lies in between; spectral
sets are pd-tiles \cite{KLMS}. Kiss, Londner, Matolcsi and Somlai ask whether every pd-tile of $\Z_M$ is a tile
and prove it for prime powers; Kadir and Fan~\cite{KF} prove it for $M=pq$ and ask about $p^nq$ and $pqr$. In
\cite{NY} it is shown, for $M=p^nq$ and for $M=pqr$, that even every \emph{weak} tile is a tile.

\begin{theorem}\label{thm:C}
Let $p\ne q$ be primes. Every weak tile of $\Z_{p^2q^2}$ is a tile.
\end{theorem}

The two proofs of \cite{NY} bound the size of a weak tile from above by a few explicit inequalities and then
count. For $p^2q^2$ that route needs a long case distinction (a mask polynomial can be divisible by any of $2^8$
sets of cyclotomic polynomials $\Phi_d$, $1<d\mid M$, and the relevant inequalities depend on whether $q<p^2$).
The proof here avoids nearly all counting.

\emph{Step 1.} If neither $\Phi_p$ nor $\Phi_q$ divides the mask polynomial of a weak tile $A\ni0$, then a weak
complement that is constant on the classes of $\gcd(x,M)$ has exactly the fraction $1/p$ of its mass on $p\Z_M$
and $1/q$ on $q\Z_M$; as $1/p+1/q<1$ it charges a unit, so no difference of two elements of $A$ is a unit, and
$A$ lies in $p\Z_M$ or in $q\Z_M$ (Lemma~\ref{lem:sands}). For tilings this is a theorem of Sands~\cite{Sa}; for
weak tilings the conclusion ``$A$ or the complement lies in a subgroup'' is false in general
(Remark~\ref{rem:sands}), but this special case survives, and it reduces the problem to $\Z_{pq^2}$ and
$\Z_{p^2q}$, where \cite{NY} applies.

\emph{Step 2.} If some weak complement of $A$ vanishes outside $p\Z_M$, the $p$ pieces of $A$ in the classes
modulo $p$ are weak tiles of $\Z_{pq^2}$ with one common weak complement. By \cite{NY} they are tiles; a common
weak complement forces them to have the same prime-power cyclotomic divisors, hence the same Coven--Meyerowitz
complement, and $A$ tiles (Lemma~\ref{lem:lift}).

\emph{Step 3.} For a weak tile with a given set $Z$ of cyclotomic divisors, the weak complements that are
constant on the nine classes are exactly the optimal points of a linear programme $LP(Z)$ in nine unknowns whose
coefficients are polynomials in $p$ and $q$, and its optimal value is $M/|A|$ (Proposition~\ref{prop:LP}). For
each of the $192$ sets $Z$ not covered by Step~1 we give an explicit dual solution and, where needed, an
explicit optimal point, as rational functions of $p$ and $q$, valid for $p<q<p^2$ or for $q>p^2$. They show one
of five things: the conditions (T1), (T2) of Coven and Meyerowitz hold; or $|A|<k$ for a number $k$ that divides
$|A|$; or there is an optimal point supported in $p\Z_M$ or in $q\Z_M$, and Step~2 applies; or (two patterns)
$|A|=pq$, and a counting lemma of Yonder gives a contradiction; or (one pattern, $q<p^2$) $|A|$ would not be an
integer.

The certificates were found by solving the programme numerically at sample primes and then solving the linear
systems of an optimal vertex symbolically. They are verified three times: by a script that checks the identities
by exact simplification and the sign conditions by a coefficient test; by exact rational arithmetic at $136$
pairs of primes; and, which makes the first two redundant, by Lean, where each certificate becomes a lemma proved
by \lean{linear\_combination} and \lean{positivity}.

\section{What is used from \texorpdfstring{\cite{NY}}{[NY]}}\label{sec:pre}

Throughout $M=p^2q^2$ and $\Z_M\cong\Z_{p^2}\times\Z_{q^2}$, $x\mapsto(u,w)$. The \emph{class} of $x$ is $(i,j)$
with $i=\min(v_p(u),2)$, $j=\min(v_q(w),2)$: nine classes, $(2,2)=\{0\}$, $(0,0)$ the units, $p\Z_M$ the union
of the classes with $i\ge1$. A frequency of class $(i,j)$ has order $p^{2-i}q^{2-j}$; we write the order
$p^\ell q^m$ as $(\ell,m)$, and say that it is a \emph{zero} of $A$ if $\hat\one_A$ vanishes at the frequencies of
that order, equivalently if $\Phi_{p^\ell q^m}$ divides $A(X)=\sum_{a\in A}X^a$. Let $Z(A)$ be the set of zeros
$(\ell,m)\ne(0,0)$,
\begin{gather*}
 E_p=\{\ell:(\ell,0)\in Z(A)\},\qquad E_q=\{m:(0,m)\in Z(A)\},\\
 F=Z(A)\cap\{1,2\}^2,\qquad k=p^{|E_p|}q^{|E_q|}.
\end{gather*}
An \emph{orphan} is an $(\ell,m)\in F$ with $\ell\notin E_p$ and $m\notin E_q$. The conditions of Coven and
Meyerowitz~\cite{CM} read: (T1) $|A|=k$; (T2) $(\ell,m)\in F$ whenever $\ell\in E_p$ and $m\in E_q$.

\begin{itemize}
\item[(P1)] If $f$ is a weak complement of $A$, then $\hat f(\xi)\hat\one_A(\xi)=0$ for $\xi\ne0$ and
$|A|\sum f=M$. Conversely, if $h\ge0$, $h(0)=1$ and $\hat h(\xi)\hat\one_A(\xi)=0$ for all $\xi\ne0$, then
$\one_A*h$ is the constant $|A|\sum h/M$.
\item[(P2)] A weak tile has a weak complement that is constant on each class (a \emph{step complement}).
\item[(P3)] $k$ divides $|A|$. A set with (T1) and (T2) tiles $\Z_M$ \cite{CM}. A weak tile without orphan
satisfies (T1).
\item[(P4)] Every weak tile of $\Z_{p^nq}$ satisfies (T1) and (T2) relative to that group and has no orphan
zero; so it tiles, and the digit set $T_{\mathcal P\setminus S_A}$ is a tiling complement, where $S_A$ is the set
of prime powers that are zeros of $A$.
\item[(P5)] For $a,b\ge1$ the order $p^aq^b$ is a zero of $A$ if and only if, with
$N_d(v)=\#\{a\in A:a\equiv v\pmod d\}$, for every $v$
\[
 pq\,N_{p^aq^b}(v)+N_{p^{a-1}q^{b-1}}(v)=q\,N_{p^{a-1}q^b}(v)+p\,N_{p^aq^{b-1}}(v).
\]
\end{itemize}
(P5) is the usual description of cyclotomic divisors by residue-class counts (in the formal proof it is derived
from the general statement \lean{Pqr.zeroAt\_iff\_cnt} of \cite{NY}). The rest is proved in \cite{NY} (Section~2 and Part~I; the absence of orphans is the lower-bound lemma of
Part~I together with (T1)), and formalised there for every cyclic group or for $p^nq$ as the case may be. We
also use the bound (S) of \cite{NY} in a form for two sets: \emph{if $g$ is a weak complement of both $A_1$ and
$A_2$, and $T$ is a nonempty set such that every common zero $d>1$ of $A_1$ and $A_2$ is a zero of $T$, then
$|A_1|\le|T|$} (the transform of $\one_T*g$ vanishes outside $0$, so it is the constant $|T|/|A_1|$, which is at
least $g(0)=1$ at a point of $T$).

\section{Two lemmas on subgroups}\label{sec:sub}

\begin{lemma}\label{lem:diff}
If $f$ is a weak complement of $A$ and $a\ne a'$ are in $A$, then $f(a-a')=0$.
\end{lemma}
\begin{proof}
$1=\sum_{a''\in A}f(a-a'')\ge f(0)+f(a-a')$.
\end{proof}

\begin{lemma}[the set lies in a subgroup]\label{lem:sands}
Let $A\ni0$ be a weak tile of $\Z_M$ such that neither $p$ nor $q$ is a zero of $A$. Then $A\subseteq p\Z_M$ or
$A\subseteq q\Z_M$, and $A$ tiles $\Z_M$.
\end{lemma}
\begin{proof}
Let $f$ be a step complement and $\sigma=\sum f$. As $\hat\one_A$ does not vanish at the order $p$, $\hat f$
does; the sum of $f$ over $p\Z_M$ is the average of $\hat f$ over the $p$ frequencies of order dividing $p$, so it
equals $\sigma/p$. Likewise the sum over $q\Z_M$ is $\sigma/q$. Hence the sum of $f$ over $p\Z_M\cup q\Z_M$ is at
most $\sigma/p+\sigma/q-f(0)<\sigma$, and $f(x)>0$ for some unit $x$, hence for all units. By
Lemma~\ref{lem:diff} no difference of two elements of $A$ is a unit. If $a\in A\setminus p\Z_M$ and
$b\in A\setminus q\Z_M$, then $q\mid a$ and $p\mid b$ because $a=a-0$ and $b$ are not units, and then $a-b$ is a
unit: a contradiction.

Say $A\subseteq p\Z_M$, and let $\iota:\Z_{M/p}\to\Z_M$, $y\mapsto py$. For $A'=\iota^{-1}(A)$ and $f'=f\circ\iota$
we get $\sum_{a'\in A'}f'(y-a')=\sum_{a\in A}f(\iota(y)-a)=1$ for all $y$: $A'$ is a weak tile of $\Z_{pq^2}$. By
(P4), $A'\oplus B'=\Z_{M/p}$, and then $A\oplus(\iota(B')+\{0,\dots,p-1\})=\Z_M$.
\end{proof}

The first paragraph is valid in $\Z_{p^aq^b}$ for all $a,b\ge1$.

\begin{lemma}[the complement lies in a subgroup]\label{lem:lift}
Let $A$ be a weak tile of $\Z_M$ with a weak complement $f$ that vanishes outside $p\Z_M$. Then $A$ tiles
$\Z_M$. The same holds with $q$ in place of $p$.
\end{lemma}
\begin{proof}
For $0\le c<p$ let $A^c=\{y\in\Z_{M/p}:c+py\in A\}$ and $g(y)=f(py)$. For $v=c+py$ the sum $\sum_{a\in A}f(v-a)$
only sees the $a\equiv c\pmod p$, so $\sum_{x\in A^c}g(y-x)=1$: every $A^c$ is a weak tile of $\Z_{pq^2}$ with the
weak complement $g$ (in particular it is not empty). By (P4) each $A^c$ satisfies (T1) and (T2) and has no
orphan; all have $n=(M/p)/\sum g$ elements. Let $S_c$ be the set of prime powers of $pq^2$ that are zeros of
$A^c$.

\emph{Claim: all $S_c$ are equal.} Let $S_1=S_c$, $S_2=S_{c'}$. By (T1) the products of the primes over $S_1$ and
over $S_2$ are both $n$; as $p$ is the only power of $p$ dividing $pq^2$, we get $p\in S_1\iff p\in S_2$. Let
$d>1$ be a zero of both pieces. There being no orphans, some $s_1\in S_1$ and some $s_2\in S_2$ divide $d$
exactly. If $s_1\in S_2$ or $s_2\in S_1$, an element of $S_1\cap S_2$ divides $d$ exactly. Otherwise neither is
$p$, so both are powers of $q$ dividing $d$ exactly, hence equal, which is absurd. So every common zero of the
two pieces is a zero of the digit set $T=T_{S_1\cap S_2}$, and by the bound (S) for two sets
$\prod_{s\in S_1\cap S_2}p_s=|T|\ge n=\prod_{s\in S_1}p_s$. Hence $S_1\subseteq S_2$, and by symmetry
$S_1=S_2$.

By (P4) the digit set $B'=T_{\mathcal P\setminus S}\subseteq\Z_{M/p}$, $S$ the common value, satisfies
$A^c\oplus B'=\Z_{M/p}$ for every $c$. Then $A\oplus pB'=\Z_M$: an element $c+py$ equals $a+pb$ with $a\in A$,
$b\in B'$ exactly when $a=c+px$ with $x\in A^c$ and $x+b=y$, which has one solution.
\end{proof}

\section{The linear programme}\label{sec:LP}

A step function is given by its values $\gs_{ij}$ on the nine classes. A character of $\Z_{p^2}$ of order
$p^\ell$, summed over the elements of valuation $i$, gives $c_p(\ell,i)$, where
\[
 c_p(0,\cdot)=(p^2-p,\ p-1,\ 1),\qquad c_p(1,\cdot)=(-p,\ p-1,\ 1),\qquad c_p(2,\cdot)=(0,\ -1,\ 1),
\]
so the Fourier transform of the step function at a frequency of order $(\ell,m)$ is
\[
 \Lambda_{\ell m}(\gs)=\sum_{i,j}\gs_{ij}\,c_p(\ell,i)\,c_q(m,j),\qquad \sigma(\gs)=\Lambda_{00}(\gs)=\sum f .
\]
For a set $Z$ of orders let $LP(Z)$ be the programme
\[
 \gs\ge0,\qquad \gs_{22}=1,\qquad \Lambda_{\ell m}(\gs)=0\ \text{ for } (\ell,m)\notin Z\cup\{(0,0)\};\qquad
 \text{minimise }\sigma(\gs).
\]

\begin{proposition}\label{prop:LP}
Let $A$ be a weak tile with $Z(A)=Z$.
\begin{enumerate}
\item Every step complement of $A$ is feasible for $LP(Z)$ and has $\sigma=M/|A|$.
\item Every feasible $\gs'$ has $\sigma(\gs')\ge M/|A|$; if $\sigma(\gs')=M/|A|$, then $\gs'$ is a step
complement of $A$.
\item \emph{(Dual certificate.)} Let $y_0$ and $y_{\ell m}$, $(\ell,m)\notin Z\cup\{(0,0)\}$, be real numbers
such that for each class $c$ the reduced cost
\[
 r_c=c_p(0,i)c_q(0,j)-y_0[c=(2,2)]-\sum_{(\ell,m)}y_{\ell m}\,c_p(\ell,i)c_q(m,j)\qquad(c=(i,j))
\]
is nonnegative. Then $\sigma(\gs)\ge y_0$ for every feasible $\gs$, so $|A|\le M/y_0$ if $y_0>0$.
\end{enumerate}
\end{proposition}
\begin{proof}
(1) is (P1). (2) Let $h$ be the step function of $\gs'$. For $\xi\ne0$ either the order of $\xi$ is in $Z$ and
$\hat\one_A(\xi)=0$, or $\hat h(\xi)=0$. By (P1), $\one_A*h$ is the constant $|A|\sigma(\gs')/M$, which at a point
of $A$ is at least $h(0)=1$. (3) $\sigma(\gs)=\sum_cr_c\gs_c+y_0\gs_{22}+\sum y_{\ell m}\Lambda_{\ell m}(\gs)
=\sum_cr_c\gs_c+y_0$.
\end{proof}

Assume $p<q$ from now on; the theorem is symmetric in the two primes. As $q\ne p^2$, exactly one of
\[
 \text{regime A: } p<q<p^2,\qquad\qquad \text{regime B: } q>p^2
\]
holds. There are $192$ sets $Z$ containing $(1,0)$ or $(0,1)$. The file \lean{certificates\_p2q2.json} gives for
each of them and each regime a \emph{category} and, except in category~a, a vector $y$ of rational functions of
$p,q$, and in the categories R and N also a point $\gs^*$:

\medskip
\begin{center}\small
\begin{tabular}{llc}
category & given, and verified for all $p,q$ in the regime & patterns (A / B)\\ \hline
a & $Z$ satisfies (T2) and has no orphan & 69 / 69\\
b & $y$ with $r\ge0$ and $y_0>M/k$ & 70 / 69\\
RP, RQ & $y$ with $r\ge0$; $\gs^*$ feasible with $\sigma(\gs^*)=y_0$ and & 22+28 / 22+30\\
 & $\gs^*=0$ outside $p\Z_M$ (RP), outside $q\Z_M$ (RQ) & \\
H & $y$ with $r\ge0$, $y_0=M/k$; $E_p=\{1\}$, $|E_q|=1$, $(2,2)\in F\not\ni(1,2)$ & 2 / 2\\
N & $y$ with $r\ge0$; $\gs^*$ feasible, $\sigma(\gs^*)=y_0=\frac{p^2(q-1)+q^2(p-1)}{pq-1}$ & 1 / 0
\end{tabular}
\end{center}
\medskip

Three examples, all in regime A. For $E_p=E_q=\{1,2\}$, $F=\emptyset$ (category b):
$y_0=p^2$, $y_{11}=1-p$, $y_{12}=0$, $y_{21}=p-q$, $y_{22}=q-p^2$; here $M/k=1$, so $|A|\le q^2<p^2q^2=k$. For
$E_p=E_q=\{1\}$, $F=\{(1,2)\}$ (category RQ; a pattern for which the programme alone allows $|A|=k$):
$y_0=pq$, $y_{02}=1-q$, $y_{20}=1-p$, $y_{22}=-(p-1)(q-1)$, the other $y$ zero, and
$\gs^*_{01}=\frac{q}{p(q-1)}$, $\gs^*_{11}=\gs^*_{21}=\frac{q-p}{p(q-1)}$, $\gs^*_{12}=\gs^*_{22}=1$, the others
zero. For $E_p=E_q=\{1\}$, $F=\{(1,2),(2,1),(2,2)\}$ (category N):
$\gs^*_{01}=\frac{q^2-p}{p(q-1)(pq-1)}$, $\gs^*_{10}=\frac{p^2-q}{q(p-1)(pq-1)}$,
$\gs^*_{11}=\frac{p^2-pq-p+q^2-q+1}{(p-1)(q-1)(pq-1)}$, $\gs^*_{22}=1$; note $\gs^*_{10}\ge0$ needs $q<p^2$.

After factoring, the sign conditions involve two dozen polynomials that are positive whenever $q>p\ge2$ (such
as $2pq-p^2-q$ or $pq^2-pq+p-q^2$), and besides them only $p^2-q$ and $2p^2q-p^2-q^2$ in regime A and $q-p^2$ in
regime B.

\section{A counting lemma}

\begin{lemma}[Yonder]\label{lem:Y}
If $|A|=pq$ and $p$ and $p^2q^2$ are zeros of $A$, then $pq^2$ is a zero of $A$.
\end{lemma}
\begin{proof}
For $c\in\Z_p$, $t\in\Z_q$ the \emph{box} $(c,t)$ is the set of $pq$ points with $u\equiv c\pmod p$,
$w\equiv t\pmod q$, a $p\times q$ array. As $p^2q^2$ is a zero, the $0/1$ array of $A$ in each box is additive by (P5), of the form
$u(c')+w(t')$; if it depended on both indices, two of its entries would differ by $2$; so it is a union of full
rows ($q$ points each) or a union of full columns ($p$ points each). As $p$ is a zero, the $q$ boxes with a given $c$ hold $|A|/p=q$ elements together. If one of
them were a union of $j\ge1$ columns and not of rows, we would get $q=pJ+qS$ with $J\ge1$, so $S=0$ and $p\mid q$.
So exactly one box $(c,t_c)$ meets $A$, in one full row. Then $A$ is invariant under the subgroup of order $q$
of $\Z_{q^2}$, and $\hat\one_A$ vanishes at every frequency that is not trivial on it, in particular at those of
order $pq^2$.
\end{proof}

\section{Proof of Theorem~\ref{thm:C}}

Let $A$ be a weak tile; translate so that $0\in A$ and let $p<q$. If neither $p$ nor $q$ is a zero,
Lemma~\ref{lem:sands} applies. Otherwise $Z(A)$ is one of the $192$ sets, and by its category:

\emph{a.} No orphan, so $|A|=k$ by (P3); with (T2), $A$ tiles by (P3).

\emph{b.} By Proposition~\ref{prop:LP}(3), $|A|\le M/y_0<k$, but $k$ divides $|A|\ge1$: there is no such weak
tile.

\emph{RP, RQ.} By Proposition~\ref{prop:LP}(3) and (1), $M/|A|\ge y_0=\sigma(\gs^*)$; by (2),
$\sigma(\gs^*)\ge M/|A|$, so $\gs^*$ is a weak complement of $A$. It vanishes outside $p\Z_M$ or $q\Z_M$, and $A$
tiles by Lemma~\ref{lem:lift}.

\emph{H.} $|A|\le M/y_0=k=pq$ and $k$ divides $|A|$, so $|A|=pq$; by Lemma~\ref{lem:Y}, $pq^2$ is a zero, but
$(1,2)\notin F$: there is no such weak tile.

\emph{N.} As in case R, $\sigma(\gs^*)=M/|A|$, so $|A|\bigl(p^2(q-1)+q^2(p-1)\bigr)=p^2q^2(pq-1)$. The number
$D=p^2(q-1)+q^2(p-1)$ is prime to $p$ and to $q$, so it divides $pq-1$; but $D-(pq-1)=(p-1)(q-1)(p+q+1)>0$:
there is no such weak tile. \qed

\section{The formal proof}\label{sec:lean}

The proof is formalised in Lean~4 (\texttt{v4.37.0-rc1}) with Mathlib, in a library \lean{P2q2} of $11{,}162$
lines built on the libraries \lean{Pqr} and \lean{Pnq} of \cite{NY}:
\begin{verbatim}
theorem P2q2.weak_tiling_iff_tiling_p2q2 (p q : ℕ)
    [Fact p.Prime] [Fact q.Prime] (hpq : p ≠ q)
    (A : Finset (ZMod (p ^ 2 * q ^ 2))) :
    (∃ f : ZMod (p ^ 2 * q ^ 2) → ℝ, Pqr.WeakTiling A f) ↔
      ∃ B : Finset (ZMod (p ^ 2 * q ^ 2)), Pqr.IsTiling A B
\end{verbatim}
with \lean{WeakTiling A f} $:=(\forall v,\,0\le f\,v)\wedge f\,0=1\wedge\forall v,\ \sum_{a\in A}f(v-a)=1$ and
\lean{IsTiling A B} $:=\forall v,\ \exists!\,(a,b),\ a\in A\wedge b\in B\wedge a+b=v$. It depends on the axioms
\lean{propext}, \lean{Classical.choice}, \lean{Quot.sound} only. A second statement in the same library (file
\lean{Pd}, $28$ lines, added on 26 November 2026),
\begin{verbatim}
theorem P2q2.tile_iff_pd_iff_weak_p2q2
\end{verbatim}
is the equivalence of tiles, pd-tiles and weak tiles of $\Z_{p^2q^2}$, where a pd-tiling complement
(\lean{Pqr.PdTiling}) is a weak tiling complement whose Fourier transform is real and nonnegative at every
frequency. It follows in a line from the theorem above and from the formalised fact that a tile is a pd-tile,
and depends on the same three axioms.

About $2{,}000$ lines are written by hand: the subgroup $p\Z_M\cong\Z_{M/p}$ and the pieces of a set
(\lean{Sub}); step functions in M\"obius coordinates, $f(x)=\sum_{d\mid\gcd(x,M)}\mu_d$, whose transform at the
order $e$ is $\sum_{e\mid d\mid M}\mu_dM/d$, for every $M$ (\lean{StepLam}), and for $M=p^2q^2$ (\lean{Nine});
(P4) and the Claim of Lemma~\ref{lem:lift} (\lean{Common}); Lemmas~\ref{lem:sands} and~\ref{lem:lift}
(\lean{Reduce}); category~a (\lean{Pattern}); Lemma~\ref{lem:Y} (\lean{LemmaY}); the arithmetic of category~N
(\lean{Arith}); the assembly (\lean{Main}). The other $9{,}132$ lines are generated from the certificate file:
$25$ polynomial inequalities (\lean{Atoms}), one lemma for each pattern and regime outside category~a ($246$ lemmas,
\lean{CertsA00}--\lean{CertsB10}), each consisting of the identity of Proposition~\ref{prop:LP}(3) with
denominators cleared, proved by \lean{linear\_combination}, and of sign conditions proved by \lean{positivity};
and the case distinction over the $256$ patterns (\lean{DispatchA}, \lean{DispatchB}). The hand-written files
other than \lean{LPDefs} and \lean{Main} were written by eight further instances, each given one file whose
statements had been fixed and compiled beforehand, and told to stop if a statement turned out false; none did.

\section{Remarks}

\begin{remark}\label{rem:sands}
For tilings $A\oplus B=\Z_M$ with $M=p^aq^b$, Sands' theorem says that $A$ or $B$ lies in a proper subgroup. The
analogue ``$A$ lies in a subgroup or every weak complement is supported in one'' fails: $A=\{0,1\}\times\{0,1,5\}
\subset\Z_4\times\Z_9$ lies in no proper subgroup, and it has step complements that charge odd numbers and
others that charge numbers prime to $3$, so the average of two of them is supported in no subgroup. What the
certificates show for $p^2q^2$ is that \emph{some} optimal complement is supported in $p\Z_M$ or $q\Z_M$ whenever
the programme does not exclude the pattern outright.
\end{remark}

\begin{remark}
In categories R, H, N and b the proof never needs to know whether the pattern is that of some set. At $M=36$,
by Yonder's census of all $2^{35}$ subsets, the $64$ patterns of weak tiles are $54$ of category~a and $10$ of
Lemma~\ref{lem:sands}.
\end{remark}

\begin{remark}
What generalises to $\Z_{p^aq^b}$: Lemma~\ref{lem:sands} as it stands; and the Claim of Lemma~\ref{lem:lift},
given (T1) and the absence of orphans for the pieces (if $S_1\ne S_2$, write $S_i=S_0\cup P_i\cup Q_i$ with
$S_0=S_1\cap S_2$ and $P_i$, $Q_i$ the powers of $p$, $q$ in $S_i\setminus S_0$; then $|P_1|=|P_2|$,
$|Q_1|=|Q_2|$, every common zero is a zero of both $T_{S_0\cup P_1\cup P_2}$ and $T_{S_0\cup Q_1\cup Q_2}$, and
one of these two sets is smaller than the pieces because $p^{|P_1|}\ne q^{|Q_1|}$). What does not: the
certificates. The computations of \cite{NY} (no weak tile violating (T1) or (T2) for $M=72,108,144,200,\dots$)
suggest that the statement holds further; a proof along these lines would need a reason, not a list, for the
existence of a complement supported in a subgroup.
\end{remark}

\begin{thebibliography}{9}
\bibitem{CM} E.~Coven and A.~Meyerowitz, \emph{Tiling the integers with translates of one finite set},
J.~Algebra 212 (1999), 161--174.
\bibitem{KF} M.~Kadir and K.~Fan, \emph{Tiles and weak tiles in $\Z_{pq}$}, arXiv:2607.02149 (2026).
\bibitem{KLMS} G.~Kiss, I.~Londner, M.~Matolcsi and G.~Somlai, \emph{Functional tilings and the
Coven--Meyerowitz tiling conditions}, arXiv:2411.03854; Discrete Analysis 2025:18.
\bibitem{KLaMS} G.~Kiss, I.~{\L}aba, C.~Marshall and G.~Somlai, \emph{Lower bounds for mask polynomials with many
cyclotomic divisors}, Adv. Math. 494 (2026), 110932; arXiv:2507.11672.
\bibitem{KMMS2} G.~Kiss, C.~Marshall, M.~Matolcsi and G.~Somlai, \emph{Nonnegative weak tilings and the
Coven--Meyerowitz condition (T2)}, arXiv:2611.13593 (2026).
\bibitem{KMSV} G.~Kiss, R.~D.~Malikiosis, G.~Somlai and M.~Vizer, \emph{On the discrete Fuglede and Pompeiu
problems}, Analysis \& PDE 13 (2020), 765--788; arXiv:1807.02844. (Abstract only.)
\bibitem{LM} N.~Lev and M.~Matolcsi, \emph{The Fuglede conjecture for convex domains is true in all
dimensions}, Acta Math. 228 (2022), 385--420.
\bibitem{NY} Noon and Yonder, \emph{Weak tiles of $\Z_{p^nq}$ and $\Z_{pqr}$ are tiles: two short proofs, checked
in Lean}, 20 November 2026 (\texttt{paper/weak-tiles-two-theorems.pdf} in this folder; the later
\texttt{paper/weak-tiles.pdf} has three theorems, this note being its Part~III).
\bibitem{Sa} A.~D.~Sands, \emph{On Keller's conjecture for certain cyclic groups}, Proc. Edinburgh Math. Soc. 22
(1979), 17--21. (Cited as it is cited in \cite{CM}; not consulted.)
\end{thebibliography}

\end{document}
